\documentclass[10pt]{article}
\usepackage[utf8]{inputenc}
\usepackage[T1]{fontenc}
\usepackage{amsmath}
\usepackage{amsfonts}
\usepackage{amssymb}
\usepackage[version=4]{mhchem}
\usepackage{stmaryrd}
\usepackage{bbold}
\usepackage{graphicx}
\usepackage[export]{adjustbox}
\graphicspath{ {./images/} }

\begin{document}
Equilibrium refinements in extensive-form games

Trembling hand perfection

\section*{Normal-Form Perfect Equilibrium}
Trembling hand in choosing a pure strategy. Fix a finite, normal-form game $\Gamma=\left\{S_{i}, u_{i}\right\}_{i=1}^{I}$.\\
Definition 1 A perturbation is a function $\epsilon: \bigcup_{i=1}^{I} S_{i} \rightarrow(0,1)$ such that $\forall i \in I, \sum_{s_{i} \in S_{i}} \epsilon\left(s_{i}\right)<1$.\\
Definition 2 Given $\Gamma$ and a perturbation $\epsilon$, define $\Gamma(\epsilon)=\left\{M_{i}(\epsilon), u_{i}\right\}_{i=1}^{I}$ where $M_{i}(\epsilon)=\left\{\mu \in M\left(S_{i}\right): \forall s_{i} \in S_{i}, \mu\left(s_{i}\right) \geq \epsilon\left(s_{i}\right)\right\}$ and $u_{i}$ is extended to $M\left(S_{i}\right)$.\\
$\Gamma(\epsilon)$ is the perturbed game with perturbation $\epsilon$.

\begin{itemize}
  \item $\Gamma(\epsilon)$ is well defined if $\forall i, M_{i}(\epsilon) \neq \emptyset$.
\end{itemize}

\section*{Example}
\begin{itemize}
  \item Consider a normal-form game with two players that we identify as Alice indexed by $\alpha$, and Bob indexed by $\beta: \Gamma=\left[S_{\alpha}, S_{\beta}, u_{\alpha}, u_{\beta}\right]$,
  \item $a$ is a generic element of $S_{\alpha}, b$ is a generic element of $S_{\beta}$.
  \item $u_{\alpha}: S_{\alpha} \times S_{\beta} \rightarrow \mathbb{R}$.
  \item Let $S_{\alpha}$ contain tree strategies, $S_{\alpha}=\left\{a_{1}, a_{2}, a_{3}\right\}$.
  \item $M\left(S_{\alpha}\right)=\left\{p=\left(p_{1}, p_{2}, p_{3}\right) \in \mathbb{R}_{+}^{3}: \sum_{i=1}^{3} p_{i}=1\right\}$ is the set of mixed strategies available to Alice.
  \item $M\left(S_{\beta}\right)=\left\{q \in \mathbb{R}_{+}^{N}: \sum_{i=1}^{N} q_{i}=1\right\}$ is the set of of mixed strategies available to Bob.
  \item Extend the payoff function to mixed strategies, $\left(\forall b \in S_{\beta}\right)\left(\forall p \in M\left(S_{\alpha}\right)\right)$,
\end{itemize}

$$
\begin{aligned}
& u_{\alpha}(p, b)=u_{\alpha}\left(a_{1}, b\right) p_{1}+u_{\alpha}\left(a_{2}, b\right) p_{2}+u_{\alpha}\left(a_{3}, b\right) p_{3} \\
& u_{\alpha}(p, b)=\sum_{i=1}^{3} u_{a}\left(a_{i}, b\right) p_{i}
\end{aligned}
$$

\begin{itemize}
  \item Similar arguments applied to Bob yield
\end{itemize}

$$
u_{\alpha}(p, q)=\sum_{a \in S_{\alpha}} \sum_{b \in S_{\beta}} u_{a}(a, b) p_{a} q_{b}
$$

\section*{Example}
\begin{center}
\includegraphics[max width=\textwidth, alt={}]{f44e9f96-45d8-4803-aeeb-6fcb841fbabe-07_1172_1859_359_153}
\end{center}

\begin{itemize}
  \item $u_{\alpha}(p, b)$ is restricted to $p$ in the smaller triangle; for $i=1,2,3$, $p_{i} \geq \epsilon\left(a_{i}\right)>0$.
\end{itemize}

$$
\begin{aligned}
& a_{1} \sim \overbrace{(1,0,0)}^{p} \rightarrow a_{1}^{\prime} \sim\left(1-\epsilon\left(a_{1}\right)-\epsilon\left(a_{2}\right), \epsilon\left(a_{2}\right), \epsilon\left(a_{3}\right)\right) \\
& a_{2} \sim(0,1,0) \rightarrow a_{2}^{\prime} \sim\left(\epsilon\left(a_{1}\right), 1-\epsilon\left(a_{1}\right)-\epsilon\left(a_{3}\right), \epsilon(3)\right) \\
& a_{3} \sim(0,0,1) \rightarrow a_{3}^{\prime} \sim\left(\epsilon\left(a_{1}\right), \epsilon\left(a_{2}\right), 1-\epsilon\left(a_{1}\right)-\epsilon\left(a_{2}\right)\right)
\end{aligned}
$$

\section*{Example. Two-player, two-strategy game}
\begin{center}
\begin{tabular}{crcc}
\multicolumn{4}{c}{$\mathbf{L}$} \\
\hline
u & 1,1 & 0,3 & $\left(1-p_{1}\right)$ \\
\hline
d & $-3,0$ & 0,0 & $\left(p_{1}\right)$ \\
\hline
 & $\left(1-q_{1}\right)$ & $q_{1}$ &  \\
\hline
\end{tabular}
\end{center}

\begin{itemize}
  \item Two NE, $(u, L)$ and $(d, R)$. Let the perturbation be
\end{itemize}

\begin{center}
\begin{tabular}{cc}
Strategy & $\epsilon(\cdot)$ \\
\hline
u & $1 / 8$ \\
\hline
d & $1 / 4$ \\
\hline
L & $1 / 16$ \\
\hline
R & $1 / 32$ \\
\hline
\end{tabular}
\end{center}

\section*{Relevant mixed strategies for $\Gamma(\epsilon)$}
\begin{center}
\includegraphics[max width=\textwidth, alt={}]{f44e9f96-45d8-4803-aeeb-6fcb841fbabe-10_1412_1412_495_150}
\end{center}

\section*{Example: Player 1's best response}
$$
\begin{aligned}
& u_{1}(p, R)=\left(1-p_{1}\right) \times 0+p_{1} \times 0 \\
& u_{1}(p, L)=\left(1-p_{1}\right) \times-3+p_{1} \times 1
\end{aligned}
$$

\begin{center}
\includegraphics[max width=\textwidth, alt={}]{f44e9f96-45d8-4803-aeeb-6fcb841fbabe-11_1033_2576_744_157}
\end{center}

\section*{Normal-form perfect equlibrium}
Definition 3 (NFPE) Let $\Gamma=\left\{S_{i}, u_{i}\right\}_{i=1}^{I}$ be a finite normal-form game. A strategy profile $\mu$ is a Normal-Form Perfect Equilibrium (NFPE) if there is a sequence of perturbations $\epsilon^{n}:\left\|\epsilon^{n}\right\|_{\infty} \rightarrow 0$ and a NE $\mu^{n}$ of $\Gamma\left(\epsilon^{n}\right)$ such that $\lim _{n \rightarrow \infty} \mu^{n}=\mu$.

\section*{Example}
\begin{center}
\begin{tabular}{lrrc}
\multicolumn{4}{c}{$\mathbf{L}$} \\
\cline { 1 - 4 }
u & 1,1 & 0,3 & $\left(1-p_{1}\right)$ \\
\hline
d & $-3,0$ & 0,0 & $\left(p_{1}\right)$ \\
\hline
 & $\left(1-q_{1}\right)$ & $q_{1}$ &  \\
\hline
\end{tabular}
\end{center}

( $d, R$ ) is NE but is not NFPE.

\begin{itemize}
  \item $(d, R)$ in $\Gamma$ is
\end{itemize}

$$
\begin{aligned}
& \bar{\mu}_{1}(u)=0=\left(1-p_{1}\right) \\
& \bar{\mu}_{2}(L)=0=\left(1-q_{1}\right)
\end{aligned}
$$

Denote $\bar{\mu}_{1}$ as $(0,1), \bar{\mu}_{2}$ as $(0,1)$.

\begin{itemize}
  \item Let $\epsilon^{n}$ be any sequence of perturbations, $\left\|\epsilon^{n}\right\|_{\infty} \rightarrow 0$.
  \item Let $\mu^{n}=\left(\mu_{1}^{n}, \mu_{2}^{n}\right)$ be strategies in $\Gamma\left(\epsilon^{n}\right), \mu_{i}^{n}\left(s_{i}\right)>\epsilon^{n}\left(s_{i}\right)>0$.
  \item Let $\mu^{n} \rightarrow \bar{\mu}=[(0,1),(0,1)]$.
\end{itemize}

$$
\begin{aligned}
E\left[U_{1} \mid u, \mu_{2}^{n}\right] & =\mu_{2}^{n}(L) \times 1+\mu_{2}^{n}(R) \times 0 \\
E\left[U_{1} \mid d, \mu_{2}^{n}\right] & =\mu_{2}^{n}(L) \times(-3)+\mu_{2}^{n}(R) \times 0
\end{aligned}
$$

\begin{itemize}
  \item $\forall \mu_{2}^{n}(L)>0, E\left[U_{1} \mid u, \mu_{2}^{n}\right]>E\left[U_{1} \mid d, \mu_{2}^{n}\right]$
  \item Then if $\mu_{1}^{n} \in B R\left(\mu_{2}^{n}\right), \mu_{1}^{n}=(1,0) \nrightarrow(0,1)$
\end{itemize}

\section*{Some properties of normal-form PE}
NFPE is commonly named Perfect Equlibrium.\\
Every NFPE is a NE.\\
Theorem 1 Every finite normal-form game has a NFPE.\\
Corollary 1 Every finite extensive-form game has a NFPE of its normalform representation.

Theorem 2 Every NE in completely mixed strategies is a NFPE.

Lemma 1 A NFPE does not use weakly dominated strategies.\\
Lemma 2 A strategy profile $\mu$ in a two person finite normal-form game is a NFPE if it is a NE and for $i=1,2, \mu_{i}$ puts no mass on any weakly dominated strategy.

\section*{Example 3-player game}
\begin{itemize}
  \item Player 1's action space is $\{T, B\}$; Player 2's is $\{L, R\}$.
  \item Player 3 chooses "the game," her action space is $\{\ell, r\}$.
\end{itemize}

\begin{center}
\begin{tabular}{llllllll}
 &  & $\mathbf{L}$ & $\mathbf{R}$ &  &  & $\mathbf{L}$ & $\mathbf{R}$ \\
\hline
$3 . \ell$ & T & $1,1,1$ & $1,0,1$ & $3 . \mathrm{r}$ & T & $1,1,0$ & $0,0,0$ \\
\hline
 & B & $1,1,1$ & $0,0,1$ &  & B & $0,1,0$ & $1,0,0$ \\
\hline
\end{tabular}
\end{center}

\begin{itemize}
  \item ( $B, L, \ell$ ) is a NE in undominated strategies; it is not NFPE.
  \item $2 . L, 3 . \ell, 1 . B$ are undominated trivially. The profile is a NE; deviating will decrease the respective players' payoffs.
  \item Exercise: Is $(T, L, \ell)$ a NE? Is it a NFPE?
\end{itemize}

\section*{( $B, L, \ell$ ) is not a NFPE}
\begin{itemize}
  \item Let $\mu_{2}^{n}, \mu_{3}^{n}$ be any pair of completely mixed strategies such that
\end{itemize}

$$
\begin{aligned}
\left(\mu_{2}^{n}(L), \mu_{2}^{n}(R)\right) & \rightarrow(1,0) & & \mathrm{L} \text { in the limit } \\
\left(\mu_{3}^{n}(\ell), \mu_{3}^{n}(r)\right) & \rightarrow(1,0) & & \ell \text { in the limit }
\end{aligned}
$$

$$
\begin{aligned}
E\left[U_{1} \mid\right. & \left.T, \mu_{2}^{n}, \mu_{3}^{n}\right]-E\left[U_{1} \mid B, \mu_{2}^{n}, \mu_{3}^{n}\right] \\
& =\sum_{s_{2}, s_{3}}\left[U_{1}\left(T, s_{2}, s_{3}\right)-U_{1}\left(B, s_{2}, s_{3}\right)\right] \mu_{2}^{n}\left(s_{2}\right) \mu_{3}^{n}\left(s_{3}\right) \\
& =(1-1) \mu_{2}^{n}(L) \mu_{3}^{n}(\ell)+(1-0) \mu_{2}^{n}(L) \mu_{3}^{n}(r) \\
& \quad+(1-0) \mu_{2}^{n}(R) \mu_{3}^{n}(\ell)+(0-1) \mu_{2}^{n}(R) \mu_{3}^{n}(r) \\
& =\mu_{2}^{n}(R) \mu_{3}^{n}(\ell)-\mu_{2}^{n}(R) \mu_{3}^{n}(r)+\mu_{2}^{n}(L) \mu_{3}^{n}(r) \\
& =\mu_{2}^{n}(R)\left[\mu_{3}^{n}(\ell)-\mu_{3}^{n}(r)\right]+\mu_{2}^{n}(L) \mu_{3}^{n}(r)>0
\end{aligned}
$$

because $\forall\left\{\left(\mu_{2}^{n}, \mu_{3}^{n}\right)\right\}_{n \geq 1}, \exists N: n>N \Longrightarrow \mu_{3}^{n}(\ell)-\mu_{3}^{n}(r)>0$.

\begin{center}
\begin{tabular}{llllllll}
 &  & $\mathbf{L}$ & $\mathbf{R}$ &  &  & $\mathbf{L}$ & $\mathbf{R}$ \\
\hline
$3 . \ell$ & T & $1,1,1$ & $1,0,1$ & $3 . \mathrm{r}$ & T & $1,1,0$ & $0,0,0$ \\
\hline
 & B & $1,1,1$ & $0,0,1$ &  & B & $0,1,0$ & $1,0,0$ \\
\hline
\end{tabular}
\end{center}

\section*{Extensive-form Perfect Equilibrium}
\begin{itemize}
  \item Trembles (mistakes) in different information sets are independent. Fix an extensive form game
\end{itemize}

$$
\Gamma=\left[\mathcal{J},(X,<),\left\{\mathcal{P}_{i}\right\}_{i=0}^{I},\left\{\mathcal{U}_{i}\right\}_{i=1}^{I},\left\{\mathcal{C}_{h}\right\}_{h \in \mathcal{U}_{i}, i \geq 1},\left\{P_{x}\right\}_{x \in \mathcal{P}_{0}}, V\right]
$$

Definition 4 A function $\epsilon: \bigcup_{h \in \bigcup \mathcal{U}_{i}} \mathcal{C}_{h} \rightarrow(0,1)$ is perturbation of $\Gamma$ if

$$
(\forall i \in I),\left(\forall h \in \mathcal{U}_{i}\right), \sum_{a \in C_{h}} \epsilon(a)<1 .
$$

Definition 5 Given an extensive-form game $\Gamma$ and a perturbation $\epsilon, \Gamma(\epsilon)$ is the extensive form game $\Gamma$ modified so that behavior strategies are limited to $\Sigma_{i}(\epsilon)=\left\{\sigma_{i} \in \Sigma_{i}: \forall h \in \mathcal{U}_{i}, \forall a \in \mathcal{C}_{h}, \sigma_{i}(a \mid h) \geq \epsilon(a)\right\}$.\\
$\Gamma(\epsilon)$ is called the perturbed game with perturbation $\epsilon$.

\begin{itemize}
  \item $\Gamma(\epsilon)$ is well defined if $\forall i, \Sigma_{i}(\epsilon) \neq \emptyset$.
\end{itemize}

\section*{Extensive-form Perfect Equlibrium}
Definition 6 (EFPE) Let $\Gamma$ be a finite extensive-form game. A behaviorstrategy profile $\sigma$ is an Extensive-Form Perfect Equilibrium (EFPE) if there is a sequence of perturbations $\epsilon^{n}$ such that $\left\|\epsilon^{n}\right\|_{\infty} \rightarrow 0$ and a NE $\sigma^{n}$ of $\Gamma\left(\epsilon^{n}\right)$ such that $\lim _{n \rightarrow \infty} \sigma^{n}=\sigma$.

Theorem 3 Every EFPE of $\Gamma$ is a SPE of $\Gamma$.\\
Theorem 4 Every finite extensive-form game with perfect recall has an extensive-form PE.

\section*{Comments}
NFPE and EFPE differ from each other.

\begin{itemize}
  \item An NFPE is a profile of distributions over pure strategies, an EFPE is a profile of functions mapping information sets to distributions over choices.
  \item An EFPE may fail to be NFPE (Exercise 7.36 in MSZ). An NFPE may fail to be EFPE (Exercise 7.37 MSZ).
  \item In EF games where each player has a single information set both NFPE and EFPE coincide.
\end{itemize}

References


\end{document}